\section{Modelagem}
\label{sec:modelagem}

\subsection{O tabuleiro como um grafo}
\label{sec:grafo}

Para tratar o problema com algoritmos conhecidos, o tabuleiro é visto como um grafo. Cada
casa é um vértice, e há uma aresta da casa A para a casa B quando o cavalo consegue ir de A
até B em um pulo. Como cada pulo conta como um movimento, todas as arestas têm o mesmo custo,
1. Assim, o menor número de movimentos de C até S é o menor número de arestas de um caminho
de C até S. Esse é o problema do \emph{menor caminho em um grafo sem pesos}.

O grafo é dirigido: como o tamanho do pulo depende do dígito de uma casa, pode haver aresta de
A para B sem haver de B para A. Também não se constrói o grafo na memória. Ele é \emph{implícito}:
dada uma casa, calculam-se na hora as casas para onde ela leva, a partir do seu dígito. Isso
evita guardar as até $8N^2$ arestas, que são 18 milhões no maior tabuleiro.

A busca em largura (BFS, do inglês \emph{breadth-first search}) resolve exatamente esse
problema~\cite{cormen}. Ela visita as casas por camadas: primeiro a casa C, a distância 0;
depois todas as casas a 1 pulo de C; depois as casas a 2 pulos; e assim por diante. Uma fila
garante essa ordem, porque as casas descobertas a partir da camada $k$ entram na fila depois
de todas as casas da camada $k$, e portanto só são visitadas depois delas. Por isso, a
primeira vez que a busca chega a S é pelo caminho mais curto: se houvesse um caminho com
menos pulos, S estaria em uma camada anterior e já teria sido encontrada.

\subsection{O tabuleiro toroidal}
\label{sec:toro}

Para dar a volta nas bordas, usa-se aritmética modular. Se o cavalo está na linha $l$ e na
coluna $c$ e faz um pulo que soma $a$ às linhas e $b$ às colunas, ele chega à linha
$(l + a) \bmod N$ e à coluna $(c + b) \bmod N$. Isso vale mesmo quando o pulo é maior que o
tabuleiro inteiro.

Por exemplo, no tabuleiro $10 \times 10$ do enunciado, uma casa com dígito 9 faz um pulo de
pernas 10 e 11. Como $10 \bmod 10 = 0$ e $11 \bmod 10 = 1$, o cavalo volta para a mesma linha
e chega à coluna seguinte: dá uma volta completa no tabuleiro e termina ao lado de onde saiu.

\subsection{A regra do pulo}
\label{sec:regra}

O enunciado não dá uma fórmula para o pulo, só a figura com os L dos dígitos 0, 1 e 2
(Figura~\ref{fig:pulos}). Por isso é preciso extrair a regra da própria figura. Nela, o L do
dígito 0 tem pernas de 2 e 1 casas, o do dígito 1 tem pernas de 3 e 2, e o do dígito 2 tem
pernas de 4 e 3. As casas coloridas são o que foi acrescentado: $d$ casas em cada perna.
Conclui-se que um dígito $d$ forma um L de pernas $1 + d$ e $2 + d$. Combinando os sinais e
trocando qual perna vai na linha e qual vai na coluna, cada casa tem oito pulos possíveis:
$(\pm(1 + d), \pm(2 + d))$ e $(\pm(2 + d), \pm(1 + d))$. Com $d = 0$ recupera-se o pulo
normal de xadrez, como o enunciado pede.

Essa leitura foi conferida no exemplo do enunciado, em que C chega a S em 3 pulos. Como o dígito
da casa C está escondido pela letra, o exemplo foi executado com cada dígito possível, de 0 a 9,
no lugar do C. Com pernas $(1 + d, 2 + d)$ o resultado é sempre 3 pulos. Também foram testadas
duas outras regras, para ver se o exemplo e a figura as distinguem (Tabela~\ref{tab:hipoteses}).

\begin{table}[ht]
  \centering
  \caption{Regras testadas para as pernas do L, com o dígito escondido sob o C variando de 0 a 9.}
  \label{tab:hipoteses}
  \small
  \begin{tabular}{|c|c|c|}
    \hline
    \textbf{Pernas para o dígito $d$} & \textbf{Pulos de C a S no exemplo} & \textbf{Compatível com a figura?}\\
    \hline
    $(1 + d,\ 2 + d)$ & 3, com todos os dígitos & sim\\
    \hline
    $(2,\ 1 + d)$ ou $(1 + d,\ 2)$ & 3 com nove dígitos, 1 com o dígito 4 & não\\
    \hline
    $(1,\ 2 + d)$ ou $(2 + d,\ 1)$ & 2, 3 ou 4, conforme o dígito & não\\
    \hline
  \end{tabular}
\end{table}

A segunda regra não é descartada pelo exemplo, mas é descartada pela figura: para $d = 1$
ela daria pernas iguais, 2 e 2, que não formam um L. Na terceira, só uma perna cresce, o que
também contraria a figura, e o resultado no exemplo dependeria do dígito escondido. Só a
primeira regra é compatível com a figura e com o exemplo para qualquer dígito.

\subsection{Duas dúvidas de interpretação}
\label{sec:duvidas}

O texto do enunciado e o exemplo não resolvem dois pontos, e para cada um é preciso escolher
uma leitura.

\textbf{De qual casa vem o dígito.} O enunciado diz que o número de cada casa afeta o quanto
o cavalo pode pular. Adota-se a leitura mais direta: o cavalo usa o dígito da casa em que está
(leitura de \emph{origem}). A leitura alternativa, em que vale o dígito da casa onde ele vai
cair (leitura de \emph{destino}), também é compatível com o enunciado e também dá 3 pulos no
exemplo. Portanto o exemplo não serve para escolher entre as duas.

\textbf{O dígito escondido.} Nos arquivos, as letras C e S ocupam o lugar de um dígito, que se
perdeu. Na leitura de origem, isso afeta o primeiro pulo, que sai da casa C; a casa S não
importa, porque o cavalo para ao chegar nela. Como o enunciado define o dígito 0 como o pulo
normal de xadrez, supõe-se o dígito 0 sob o C, isto é, o primeiro pulo é um pulo normal de
cavalo. Na leitura de destino, o dígito escondido que importa é o de S, e faz-se a mesma
suposição.

\textbf{Quanto isso importa.} A Tabela~\ref{tab:sensibilidade} mostra o resultado de cada caso
nas quatro combinações: leitura de origem ou de destino, com o dígito escondido valendo 0 ou
sendo qualquer dígito de 0 a 9 (nesse caso, toma-se o menor resultado entre os dez). Só o
caso 1500 dá o mesmo resultado nas quatro combinações; nos outros sete, a resposta depende da
escolha.

\begin{table}[ht]
  \centering
  \caption{Menor número de movimentos de cada caso, segundo a leitura do enunciado e o dígito
  escondido. A primeira coluna é a leitura adotada neste relatório.}
  \label{tab:sensibilidade}
  \small
  \begin{tabular}{|c|c|c|c|c|}
    \hline
    & \multicolumn{2}{c|}{\textbf{Leitura de origem}} & \multicolumn{2}{c|}{\textbf{Leitura de destino}}\\
    \cline{2-5}
    \textbf{Caso} & \textbf{C = 0 (adotada)} & \textbf{C qualquer} & \textbf{S = 0} & \textbf{S qualquer}\\
    \hline
    40 & 4 & 4 & 2 & 2\\
    \hline
    80 & 6 & 4 & 6 & 4\\
    \hline
    100 & 8 & 6 & 8 & 6\\
    \hline
    150 & 11 & 9 & 11 & 11\\
    \hline
    200 & 9 & 9 & 11 & 9\\
    \hline
    400 & 19 & 17 & 17 & 17\\
    \hline
    800 & 43 & 43 & 43 & 41\\
    \hline
    1500 & 34 & 34 & 34 & 34\\
    \hline
  \end{tabular}
\end{table}

Por isso os números deste relatório são apresentados como os da leitura de origem com o dígito 0,
e a dependência fica registrada. Se o enunciado for esclarecido, o algoritmo da
Seção~\ref{sec:algoritmo} não muda: muda apenas a regra que diz para onde cada casa leva.
\pendente{Confirmar com o professor qual leitura vale. Se mudar, ajustar esta seção, o Resumo
e a Tabela~\ref{tab:resultados}.}
